\section{Conclusion}

We set out to test whether machine learning research suffers from epistemic lock-in---a self-reinforcing cycle where benchmark concentration compounds over time. Our findings reveal a more nuanced picture: the field co-evolves with its evaluation standards without becoming trapped by them.

Our analysis of 12,600 papers across 21 venue-years yields three principal contributions:

\begin{enumerate}
    \item \textbf{Benchmark concentration is measurable and systematic.} HHI provides a tractable metric, revealing moderate concentration (range 0.007--0.046) with high correlation to top-5 dataset dominance ($\rho$=0.90).

    \item \textbf{Citation networks serve as conduits for benchmark propagation.} Papers that cite each other exhibit 24-fold higher benchmark overlap (Jaccard 0.318 vs.\ 0.014, Cohen's $d$=1.93).

    \item \textbf{No evidence for temporal lock-in.} The lagged regression coefficient is positive rather than negative ($\beta$=+1.603, $p$=0.098). Granger causality tests show zero significant venues.
\end{enumerate}

The fear of epistemic lock-in assumes that benchmark adoption is a one-way ratchet. Our evidence suggests otherwise: the machine learning community collectively selects evaluation standards, revises them, and moves on---co-evolving with its benchmarks rather than being captured by them.
